%% GENERATED by tools/make_aggregation_robustness.py -- do not hand-edit.

\section{Does the headline depend on how the ratio was averaged?}
%% No \label here: both including documents already own a label for the
%% surrounding chapter/appendix, and defining one in the shared fragment made
%% it multiply-defined in the thesis.
\label{sec:aggregation-headline}

A ratio is a nonlinear function of two means, so its value depends on the order
of averaging. \headlinetable{} reports the schedule-pooled paired
per-seed ratio: squared score error and reference score energy are summed over
the twelve noise levels before the square root, the ratio is formed within each
seed, and the seed is the inferential unit --- the twelve levels share a
training set and a fitted model. Being the definition stated in the text does
not make it the only defensible summary, so here is the same data under four
alternatives, with a paired interval on the one the table reports.

\begin{center}
%% minipage, not a bare center: without it LaTeX will break between the
%% caption and the tabular, and it did -- Table 9.5's caption sat alone at
%% the foot of one page with its rows at the head of the next.
\begin{minipage}{\linewidth}\small\centering
\setlength{\tabcolsep}{3pt}
\captionof{table}[The headline ratio under five estimands]{The headline under five estimands. (0) is what
\headlinetable{} reports; (1) is the mean of per-(seed, level)
ratios, the summary earlier drafts reported. The interval is a paired
bootstrap over seeds, 20,000 resamples; seeds are resampled rather than
cells, because cells within a seed share a fitted model.}
\label{tab:aggregation}
\begin{tabular}{rcccccc}
\toprule
$\nseq$ & (0) pooled & (1) per-cell & (2) ratio of means & (3) geometric & (4) median
        & (5) 95\% CI on (0) \\
\midrule
32 & $6.9$ & $7.0$ & $6.8$ & $6.1$ & $6.8$ & $[6.4,\,7.5]$ \\
64 & $6.6$ & $7.8$ & $6.6$ & $6.7$ & $6.6$ & $[5.6,\,7.6]$ \\
128 & $7.9$ & $8.8$ & $8.0$ & $7.8$ & $8.6$ & $[7.0,\,8.8]$ \\
256 & $8.2$ & $9.7$ & $8.8$ & $8.6$ & $9.2$ & $[7.3,\,9.0]$ \\
512 & $8.2$ & $9.0$ & $8.6$ & $7.9$ & $8.2$ & $[7.2,\,9.2]$ \\
1024 & $9.0$ & $12.5$ & $9.7$ & $10.8$ & $11.1$ & $[7.7,\,10.4]$ \\
2048 & $8.6$ & $15.5$ & $10.0$ & $13.2$ & $12.9$ & $[7.4,\,9.9]$ \\
4096 & $10.9$ & $17.6$ & $12.4$ & $15.6$ & $16.0$ & $[9.4,\,12.5]$ \\
\bottomrule
\end{tabular}
\end{minipage}
\end{center}

\noindent Two things follow.

The conclusion does not depend on the choice: no estimand at any size falls
below $6.1$, every bootstrap interval sits well clear of $1$, and no
aggregation produces a reversal anywhere.

But the estimands are far from interchangeable, and the mean of cellwise
ratios (1) is the largest of the five at every size, reaching
$17.6$ at the largest size where the pooled estimand reads
$\ratiohi$. The reason is mechanical: a per-level ratio at a level with a
small reference score norm can be very large while contributing little
energy, and averaging ratios weights all levels equally. Pooling the
energies first weights each level by how much score there is to get wrong,
which is exactly what the pooled definition specifies. Earlier drafts reported
(1) as the headline; the range $7.0$--$17.6$ they quoted is
therefore an artefact of that estimand, and the pooled range
$\ratiolo$--$\ratiohi$ replaces it. The ratio of per-seed averaged
errors (2) runs $6.6$ to $12.4$.

No inequality orders these summaries in general: Jensen gives
$\E[1/Y]\ge 1/\E[Y]$, but $\E[X/Y]\ge\E[X]/\E[Y]$ does \emph{not}
follow for arbitrary positive correlated $X$ and $Y$. The observed ordering
is a property of these cells, which is exactly why the estimand has to be
named rather than assumed.

Resolving the average over noise levels shows where any single number hides
structure: the per-level ratio is not flat in $t$ but peaks near
$t \approx 0.5$ and falls at both ends. Two different ``weakest''
quantities have to be kept apart. The weakest seed-averaged $(\nseq,t)$
aggregate in the grid is $3.6$; the weakest \emph{individual}
cell is $0.92$, and EM--BP is beaten outright in
$2$ of the $1,520$ cells summarised here. The estimator's
advantage is largest at moderate noise, where the posterior is dominated by
neither the likelihood nor the prior --- which is where knowing the
transition structure should help most.
